\section{Introduction}\label{sec:intro}

\paragraph{Problem.}
Diagnostic models for skin lesions are typically trained on curated dermoscopic images, yet real
consumer-device photographs span a wide quality range --- motion blur, poor illumination, incomplete
framing, colour casts, and low contrast --- on which a model accurate on clean data can be both wrong
and confident. The deployment question is therefore not only ``what is the diagnosis?'' but ``is this
image good enough to diagnose, and if not, what should the system do about it?''

\paragraph{Why existing tools fall short.}
Three families each address part of the problem but none closes the loop:
(i) \emph{post-hoc calibration} (temperature scaling \citep{guo2017calibration}, Platt/isotonic) and
(ii) \emph{uncertainty estimation} (MC dropout \citep{gal2016dropout}, deep ensembles
\citep{lakshminarayanan2017simple}) both yield quality-oblivious confidence, unable to attribute it to
degradation rather than genuine ambiguity; (iii) \emph{generic enhancement} (Real-ESRGAN
\citep{wang2021realesrgan}, Restormer \citep{zamir2022restormer}) optimises perception with no
guarantee that diagnostic content survives, and generative restoration risks inventing clinically
unacceptable structure (pigment networks, vessels). None \emph{preserves the diagnostic signal under a
guarantee}; gradient-level repair~\citep{gradprom2025} treats the symptom during training, whereas we
bound the information loss directly and decide, per input, when enhancement should be withheld in
favour of re-acquisition.

\paragraph{Closed-loop quality-triage.}
We argue that quality must be a first-class signal threaded through the whole pipeline. Our system
forms a closed loop: \visiscore assesses quality; the agent decides whether to diagnose directly,
diagnose with caution, enhance-then-diagnose with our diagnosis-preserving enhancer \visienhance, or
query for a retake; and a frozen classifier produces the final diagnosis, its posterior confidence
being the signal the gate routes on (a standard bottleneck posterior suffices, \S\ref{sec:exp-dca}). Crucially, the driving quality scalar
need not come from \visiscore alone: any per-input estimate --- a learned scorer, a no-reference IQA
metric, or a known corruption severity --- keeps the framework general.

\paragraph{Contributions.}
To motivate quality-aware triage, we first chart a reliability~$\times$~recoverability decision
surface (per-dimension~$\times$~per-severity, AUC-led, recoverability axis derived from enhancement)
that quantifies where degradation breaks diagnosis and where enhancement recovers it; complementing
rather than reusing prior per-dimension calibration views, it frames the three contributions below.
\begin{itemize}
\item \textbf{C1 --- A diagnosis-preserving enhancement objective with a mutual-information lower
bound.} A feature-level objective (DP-Loss) constrains the enhancer \visienhance to retain the
information a frozen diagnostic probe relies on, rather than optimising perceptual fidelity. We derive
a Pinsker-optimal mutual-information lower bound, $I(\Zenh;Y)\geq I(\Zref;Y)-\beta\sqrt{\epsilon}$
under a KL budget $\Ldp\leq\epsilon$ (Lemma~\ref{lem:dp}), with a directional entropy result
(Proposition~\ref{prop:entropy}). An ablation attributes the effect to the objective itself
($\Delta\mathrm{AUC}_{\text{enh}}$ improves by $+0.0205$, $95\%$ CI $[+0.005,+0.035]$; McNemar
$p=2.3\times10^{-45}$), and against six restoration baselines under a matched zero-shot protocol
\visienhance preserves diagnosis better on every one\todo{R1 核：fine-tuned baseline 公平比数字补强后升级此句口径}.
\item \textbf{C2 --- A query-for-retake agent.} A routing policy selects among four channels ---
direct diagnosis, cautioned diagnosis, enhance-then-diagnose, and query-for-retake --- of which
query-for-retake is, to our knowledge, the first to request \emph{re-acquisition} rather than merely
abstain or defer. A local, per-band conditional risk guarantee (Theorem~\ref{thm:agent-compact})
backs its licence to enhance: the benefit is positive only inside a quality band
$[\tauenh,\tauhigh]$, outside which the policy routes to a retake.
\item \textbf{C3 --- An honest reliability boundary.} We show, and do not hide, where enhancement
turns net-negative: severely degraded malignant cases, where reweighting the reconstruction loss
fails to close the salvage gap (a melanoma salvage of $5.2\%$ against a $31\%$ damage rate), and
extreme contrast loss at the per-dimension level. These boundaries are the hardest motivation for the
query-for-retake gate; correspondingly, under a fixed-threshold policy global triage does not yet beat
direct diagnosis (confidence-gated sensitivity $0.818$ versus $0.788$), which we report as motivation
for learned routing thresholds rather than a result to suppress.
\end{itemize}
